\documentclass[a4paper,12pt,twoside]{memoir}

% Castellano
\usepackage[spanish,es-tabla]{babel}
\selectlanguage{spanish}
\usepackage[utf8]{inputenc}
\usepackage[T1]{fontenc}
\usepackage{lmodern} % scalable font
\usepackage{microtype}
\usepackage{placeins}

\RequirePackage{booktabs}
\RequirePackage[table]{xcolor}
\RequirePackage{xtab}
\RequirePackage{multirow}

% Links
\PassOptionsToPackage{hyphens}{url}\usepackage[colorlinks]{hyperref}
\hypersetup{
	allcolors = {red}
}

% Ecuaciones
\usepackage{amsmath}

% Rutas de fichero / paquete
\newcommand{\ruta}[1]{{\sffamily #1}}

% Párrafos
\nonzeroparskip

% Huérfanas y viudas
\widowpenalty100000
\clubpenalty100000

% Evitar solapes en el header
\nouppercaseheads

% Imagenes
\usepackage{graphicx}
\newcommand{\imagen}[3]{
	\begin{figure}[!h]
		\centering
		\includegraphics[width=#3\textwidth]{#1}
		\caption{#2}\label{fig:#1}
	\end{figure}
	\FloatBarrier
}

\newcommand{\imagenflotante}[2]{
	\begin{figure}%[!h]
		\centering
		\includegraphics[width=0.9\textwidth]{#1}
		\caption{#2}\label{fig:#1}
	\end{figure}
}



% El comando \figura nos permite insertar figuras comodamente, y utilizando
% siempre el mismo formato. Los parametros son:
% 1 -> Porcentaje del ancho de página que ocupará la figura (de 0 a 1)
% 2 --> Fichero de la imagen
% 3 --> Texto a pie de imagen
% 4 --> Etiqueta (label) para referencias
% 5 --> Opciones que queramos pasarle al \includegraphics
% 6 --> Opciones de posicionamiento a pasarle a \begin{figure}
\newcommand{\figuraConPosicion}[6]{%
  \setlength{\anchoFloat}{#1\textwidth}%
  \addtolength{\anchoFloat}{-4\fboxsep}%
  \setlength{\anchoFigura}{\anchoFloat}%
  \begin{figure}[#6]
    \begin{center}%
      \Ovalbox{%
        \begin{minipage}{\anchoFloat}%
          \begin{center}%
            \includegraphics[width=\anchoFigura,#5]{#2}%
            \caption{#3}%
            \label{#4}%
          \end{center}%
        \end{minipage}
      }%
    \end{center}%
  \end{figure}%
}

%
% Comando para incluir imágenes en formato apaisado (sin marco).
\newcommand{\figuraApaisadaSinMarco}[5]{%
  \begin{figure}%
    \begin{center}%
    \includegraphics[angle=90,height=#1\textheight,#5]{#2}%
    \caption{#3}%
    \label{#4}%
    \end{center}%
  \end{figure}%
}
% Para las tablas
\newcommand{\otoprule}{\midrule [\heavyrulewidth]}
%
% Nuevo comando para tablas pequeñas (menos de una página).
\newcommand{\tablaSmall}[5]{%
 \begin{table}
  \begin{center}
   \rowcolors {2}{gray!35}{}
   \begin{tabular}{#2}
    \toprule
    #4
    \otoprule
    #5
    \bottomrule
   \end{tabular}
   \caption{#1}
   \label{tabla:#3}
  \end{center}
 \end{table}
}

%
%Para el float H de tablaSmallSinColores
\usepackage{float}

%
% Nuevo comando para tablas pequeñas (menos de una página).
\newcommand{\tablaSmallSinColores}[5]{%
 \begin{table}[H]
  \begin{center}
   \begin{tabular}{#2}
    \toprule
    #4
    \otoprule
    #5
    \bottomrule
   \end{tabular}
   \caption{#1}
   \label{tabla:#3}
  \end{center}
 \end{table}
}

\newcommand{\tablaApaisadaSmall}[5]{%
\begin{landscape}
  \begin{table}
   \begin{center}
    \rowcolors {2}{gray!35}{}
    \begin{tabular}{#2}
     \toprule
     #4
     \otoprule
     #5
     \bottomrule
    \end{tabular}
    \caption{#1}
    \label{tabla:#3}
   \end{center}
  \end{table}
\end{landscape}
}

%
% Nuevo comando para tablas grandes con cabecera y filas alternas coloreadas en gris.
\newcommand{\tabla}[6]{%
  \begin{center}
    \tablefirsthead{
      \toprule
      #5
      \otoprule
    }
    \tablehead{
      \multicolumn{#3}{l}{\small\sl continúa desde la página anterior}\\
      \toprule
      #5
      \otoprule
    }
    \tabletail{
      \hline
      \multicolumn{#3}{r}{\small\sl continúa en la página siguiente}\\
    }
    \tablelasttail{
      \hline
    }
    \bottomcaption{#1}
    \rowcolors {2}{gray!35}{}
    \begin{xtabular}{#2}
      #6
      \bottomrule
    \end{xtabular}
    \label{tabla:#4}
  \end{center}
}

%
% Nuevo comando para tablas grandes con cabecera.
\newcommand{\tablaSinColores}[6]{%
  \begin{center}
    \tablefirsthead{
      \toprule
      #5
      \otoprule
    }
    \tablehead{
      \multicolumn{#3}{l}{\small\sl continúa desde la página anterior}\\
      \toprule
      #5
      \otoprule
    }
    \tabletail{
      \hline
      \multicolumn{#3}{r}{\small\sl continúa en la página siguiente}\\
    }
    \tablelasttail{
      \hline
    }
    \bottomcaption{#1}
    \begin{xtabular}{#2}
      #6
      \bottomrule
    \end{xtabular}
    \label{tabla:#4}
  \end{center}
}

%
% Nuevo comando para tablas grandes sin cabecera.
\newcommand{\tablaSinCabecera}[5]{%
  \begin{center}
    \tablefirsthead{
      \toprule
    }
    \tablehead{
      \multicolumn{#3}{l}{\small\sl continúa desde la página anterior}\\
      \hline
    }
    \tabletail{
      \hline
      \multicolumn{#3}{r}{\small\sl continúa en la página siguiente}\\
    }
    \tablelasttail{
      \hline
    }
    \bottomcaption{#1}
  \begin{xtabular}{#2}
    #5
   \bottomrule
  \end{xtabular}
  \label{tabla:#4}
  \end{center}
}



\definecolor{cgoLight}{HTML}{EEEEEE}
\definecolor{cgoExtralight}{HTML}{FFFFFF}

%
% Nuevo comando para tablas grandes sin cabecera.
\newcommand{\tablaSinCabeceraConBandas}[5]{%
  \begin{center}
    \tablefirsthead{
      \toprule
    }
    \tablehead{
      \multicolumn{#3}{l}{\small\sl continúa desde la página anterior}\\
      \hline
    }
    \tabletail{
      \hline
      \multicolumn{#3}{r}{\small\sl continúa en la página siguiente}\\
    }
    \tablelasttail{
      \hline
    }
    \bottomcaption{#1}
    \rowcolors[]{1}{cgoExtralight}{cgoLight}

  \begin{xtabular}{#2}
    #5
   \bottomrule
  \end{xtabular}
  \label{tabla:#4}
  \end{center}
}




\graphicspath{ {./img/} }

% Capítulos
\chapterstyle{bianchi}
\newcommand{\capitulo}[2]{
	\setcounter{chapter}{#1}
	\setcounter{section}{0}
	\setcounter{figure}{0}
	\setcounter{table}{0}
	\chapter*{#2}
	\addcontentsline{toc}{chapter}{#2}
	\markboth{#2}{#2}
}

% Apéndices
\renewcommand{\appendixname}{Apéndice}
\renewcommand*\cftappendixname{\appendixname}

\newcommand{\apendice}[1]{
	%\renewcommand{\thechapter}{A}
	\chapter{#1}
}

\renewcommand*\cftappendixname{\appendixname\ }

% Formato de portada
\makeatletter
\usepackage{xcolor}
\newcommand{\tutor}[1]{\def\@tutor{#1}}
\newcommand{\cotutor}[1]{\def\@cotutor{#1}}
\newcommand{\course}[1]{\def\@course{#1}}
\definecolor{cpardoBox}{HTML}{E6E6FF}
\def\maketitle{
  \null
  \thispagestyle{empty}
  % Cabecera ----------------
\noindent\includegraphics[width=\textwidth]{cabecera}\vspace{1cm}%
  \vfill
  % Título proyecto y escudo informática ----------------
  \colorbox{cpardoBox}{%
    \begin{minipage}{.8\textwidth}
      \vspace{.5cm}\Large
      \begin{center}
      \textbf{TFG del Grado en Ingeniería Informática}\vspace{.6cm}\\
      \textbf{\LARGE\@title{}}
      \end{center}
      \vspace{.2cm}
    \end{minipage}

  }%
  \hfill\begin{minipage}{.20\textwidth}
    \includegraphics[width=\textwidth]{escudoInfor}
  \end{minipage}
  \vfill
  % Datos de alumno, curso y tutores ------------------
  \begin{center}%
  {%
    \noindent\LARGE
    Presentado por \@author{}\\ 
    en Universidad de Burgos --- \@date{}\\
    Tutor: \@tutor{}\\
    Cotutor: \@cotutor{}\\
  }%
  \end{center}%
  \null
  \cleardoublepage
  }
\makeatother


% Datos de portada
\title{Digitalización del Boletín de Turismo del Observatorio de la Provincia de Burgos \\Documentación Técnica}
\author{Christian Andrés Núñez Duque}
\tutor{Bruno Baruque Zanón}
\cotutor{Julio César Puche Regaliza}
\date{\today}

\begin{document}

\maketitle



\cleardoublepage



%%%%%%%%%%%%%%%%%%%%%%%%%%%%%%%%%%%%%%%%%%%%%%%%%%%%%%%%%%%%%%%%%%%%%%%%%%%%%%%%%%%%%%%%



\frontmatter


\clearpage

% Indices
\tableofcontents

\clearpage

\listoffigures

\clearpage

\listoftables

\clearpage

\mainmatter

\appendix

\apendice{Plan de Proyecto Software}

\section{Introducción}

En este apéndice se pretende detallar el plan de proyecto software seguido para el desarrollo del mismo. 
Para ello, se mostrará la planificación temporal, el estudio de viabilidad que comprende aspectos relevantes de su viabilidad económica y legal.

\section{Planificación temporal}

El proyecto se ha desarrollado siguiendo el modelo de desarrollo ágil SCRUM. Para ello, se han definido una serie de sprints de 2 semanas utilizando GitHub Projects.
En cada sprint se ha tratado de implementar una serie de funcionalidades definidas previamente como issues.

\subsection{Sprint 1 (27/02/2025 - 12/03/2025)}

Este primer sprint se centró en la creación del proyecto y en su inicio. Dado que este proyecto se basa en un TFG anterior, algunas tareas están relacionadas con ellas.

\begin{itemize}
    \item Leer y analizar el TFG previo relacionado con el tema.
    \item Definir el enfoque y tecnologías a utilizar en el proyecto.
    \item Probar diferentes APIs (de extracción de datos).
    \item Probar funcionalidades de Power BI.
    \item Probar extracciones de reviews a partir de GeoJSON de OpenStreetMaps.
    \item Probar funcionalidades de PowerPages.
    \item Probar la viabilidad de extraer datos con Apify.
    \item Probar Google Maps Scraper.
\end{itemize}

\subsection{Sprint 2 (13/03/2025 - 26/03/2025)}

En el segundo sprint, se trató de iniciar el proceso de documentación del proyecto y de iniciar el sistema de extracción de datos.

\begin{itemize}
    \item Obtener place ids sin usar los cuadrantes de las coordenadas.
    \item Iniciar el proceso de documentación
    \item Obtener place ids a partir de textSearch.
    \item Documentar los objetivos del proyecto.
    \item Actualizar documentación de los objetivos del proyecto.
    \item Automatizar la extracción de los datos del INE a través de su API. (Cerrada como no planeada)
\end{itemize}

\subsection{Sprint 3 (27/03/2025 - 09/04/2025)}

Para el tercer sprint, el objetivo principal era diseñar el modelo de datos, el cual fue actualizado en los últimos sprints.

\begin{itemize}
    \item Diseñar un esquema para el modelo de datos.
    \item Añadir manualmente los datos (poblacion y coordenadas) de los municipios que no provee el INE. (Cerrada como no planeada)
\end{itemize}

\subsection{Sprint 4 (10/04/2025 - 23/04/2025)}

En el cuarto sprint, se inició a desarrollar el sistema de aprendizaje automático y el primer modelo de red neuronal.

\begin{itemize}
    \item Crear una red neuronal a partir de reviews.
\end{itemize}

\subsection{Sprint 5 (24/04/2025 - 07/05/2025)}

En el quinto sprint se trató de crear el modelo final de extracción de datos.

\begin{itemize}
    \item Probar la extracción de reviews mediante Overpass QL, Google Places API y Apify.
    \item Crear una base de datos para almacenar las reviews.
\end{itemize}

\subsection{Sprint 6 (08/05/2025 - 21/05/2025)}

Durante el sexto sprint, tras no haber obtenido buenos resultados con el modelo de red neuronal, se optó por construir un dataset propio para el modelo de aprendizaje automático.

\begin{itemize}
    \item Crear un dataset con reviews y categorías de los POIs.
\end{itemize}

\subsection{Sprint 7 (22/05/2025 - 04/06/2025)}

El séptimo sprint se basó en la creación de un diccionario para mapear subcategorías a categorías para mejorar el modelo de aprendizaje automático.
Finalmente esto fue descartado tras rediseñar el conjunto de datos.

\begin{itemize}
    \item Hacer un mapeo de categorías a subcategorías (Cerrada como no planeada).
\end{itemize}

\subsection{Sprint 8 (05/06/2025 - 18/06/2025)}

El octavo sprint se centró en redactar la documentación de la memoria del proyecto. El objetivo era crear una versión inicial pero completa de la memoria.

\begin{itemize}
    \item Añadir conceptos teóricos a la documentación.
    \item Añadir técnicas y herramientas a la documentación.
    \item Añadir aspectos relevantes a la documentación.
    \item Añadir los trabajos relacionados a la documentación.
\end{itemize}

\subsection{Sprint 9 (19/06/2025 - 02/07/2025)}

El penúltimo sprint se centró en completar la memoria y crear el cuadro de mando.

\begin{itemize}
    \item Añadir la introducción a la documentación.
    \item Añadir las conclusiones y líneas de trabajo futuras a la documentación.
    \item Crear un dashboard de Power BI para la visualización de resultados.
    \item Crear una Azure Function para la predicción y guardados automáticos sobre los recursos.
\end{itemize}

\subsection{Sprint 10 (03/07/2025 - 07/07/2025)}

En el último sprint se ha tratado de completar toda la documentación y el cuadro de mando.
Además se ha ordenado el repositorio y se han actualizado los ficheros antiguos.

\begin{itemize}
    \item Implementar un sistema de análisis y detección de defectos de código.
    \item Crear un sistema de recomendación a partir de la localización y reseñas de los recursos.
    \item Añadir el plan de proyecto a los anexos.
    \item Añadir los requisitos a los anexos.
    \item Añadir la especificación de diseño a los anexos.
    \item Añadir el manual de programador a los anexos.
    \item Añadir el manual de usuario a los anexos.
    \item Añadir el anexo de sostenibilización curricular.
\end{itemize}

\imagen{commits}{Gráfica de commits}{1}

\section{Estudio de viabilidad}

\subsection{Viabilidad económica}

Desde el punto de vista económico, hay varios aspectos a tener en cuenta:

\subsubsection{Costes de hardware}

Este proyecto ha sido desarrollado en un ordenador de sobremesa con un coste aproximado de 800 euros comprado hace 9 años por lo que se puede considerar amortizado.

\subsubsection{Extracción de datos}

Para la extracción de datos se ha optado por utilizar la plataforma Apify que permite extraer datos de forma masiva de páginas web y APIs.
Se muestra una tabla con los precios de las cuentas de Apify\cite{apify:pricing} en dólares américanos:
\begin{table}[H]
    \centering
    \begin{tabular}{|c|c|c|c|c|}
        \hline
        \textbf{Tipo de cuenta} & \textbf{Número de reviews} & \textbf{Precio/mes al 50\%} \\
        \hline
        Free & 14285 & 0\$ \\
        Starter & 111428 & 19,5\$ \\
        Scale & 568571 & 99,5\$ \\
        Business & 2854285 & 499,5\$ \\
        Enterprise & Ilimitado & Hablar con Apify \\
        \hline
    \end{tabular}
    \caption{Costos por solicitud en Apify}
\end{table}
Cabe destacar que ofrecen un descuento del 50\% para cuentas educativas.

Se estiman alrededor de un millón y medio de reseñas en toda la provincia de Burgos, por lo que la mejor opción puede parecer la cuenta Business.
Sin embargo, hay muchos recursos con pocas reseñas o sin reseñas que podrían ser descartados por lo que la mejor opción serían 2 cuentas Scale.

Estos pagos son mensuales, por lo que se estima un coste de 198\$ al mes en caso de que cada mes se quiera realizar una nueva extracción.
Haciendo la conversión a euros, a día 3 de julio de 2025, el coste sería de \textbf{167,92 euros al mes}. (1 USD = 0,85 EUR)

Si se decide pagar anualmente, se aplicaría un descuento del 10\% sobre cada cuenta, lo que supondría un coste de 179\$ al mes entre las dos cuentas y de 2148\$ anualmente.
Haciendo la conversión a euros, el coste sería de \textbf{1821,63 euros al año}. (1 USD = 0,85 EUR) 

Estos costes han sido calculados teniendo en cuenta que se aplicaría un descuento del 50\% por cuentas educativas.
Cabe destacar que estos cálculos están pensados para una futura escalabilidad del proyecto. Para el desarrollo del mismo se ha optado por utilizar cuentas gratuitas y un scraper de código abierto gratuito.

\subsubsection{Costes de alojamiento}

Para el alojamiento de la base de datos, se ha optado por utilizar un servidor en la nube de Microsoft Azure.
La opción elegida ha sido la opción serverless ya que es la más económica y permite alojar la base de datos correctamente.
En el modelo Serverless, se paga 5,69 USD al mes por 41,6 GB de almacenamiento, además de un coste adicional de 0,000159 USD por cada segundo de procesamiento de la CPU. En cambio, con un servidor dedicado, el coste mensual asciende a 396,10 USD por los mismos 41,6 GB de almacenamiento, incluyendo dos núcleos de CPU.
Respecto al alojamiento del cuadro de mando en PowerPages es necesaria una cuenta de Power BI Pro que tiene un coste de 13,10 euros al mes.
De esta forma se estima un coste mensual de entre 6 y 8 USD, lo que al cambio a euros sería aproximadamente de \textbf{5 a 7 euros al mes}. (1 USD = 0,85 EUR)

\subsection{Viabilidad legal}

\subsubsection{Protección de datos}

Respecto a la protección de datos, se ha procurado mantener en el anónimato los datos extraídos. 
No se trata de un proyecto con ánimo de lucro por lo que no se pretende obtener ningún beneficio económico de los datos extraídos.
En ningún momento del cuadro de mando se muestra información personal de los usuarios, únicamente se muestra la distinción entre género.

\subsubsection{Extracción de datos}

Si bien es cierto que la extracción de datos mediante scraper es un tema controvertido, en España, realizar web scraping es completamente legal.
En este proyecto, la información extraída es de acceso público y no se está vulnerando ningún derecho de propiedad intelectual.\cite{webscraping}

\subsubsection{Licencias de software}

Durante el desarrollo de este proyecto se han utilizado las siguientes herramientas con sus correspondientes licencias software:

\begin{table}[H]
	\centering
	\begin{tabularx}{\linewidth}{ p{0.4\columnwidth} p{0.3\columnwidth} p{0.3\columnwidth} }
		\toprule
		\textbf{Herramienta}    & \textbf{Versión}      &\textbf{Licencia}\\
        \toprule
		Visual Studio Code      & 1.101.2               & MIT \\
        LaTeX                   & 2023                  & LPPL \\
        Google Colab            & N/A                   & Propietaria \\
        Git                     & 2.47.1                & GPLv2 \\
        GitHub                  & N/A                   & Propietaria \\
        Python                  & 3.10.11               & PSF \\
        TensorFlow              & 2.18.0                & Apache 2.0 \\
        Keras                   & 3.8.0                 & Apache 2.0 \\
        Numpy                   & 2.0.2                 & BSD \\
        Scikit Learn            & 1.6.1                 & BSD \\
        Pysentimiento           & 0.7.3                 & MIT \\
        Emoji                   & 2.14.1                & BSD \\
        Pyodbc                  & 5.2.0                 & MIT \\
        Langdetect              & 1.0.9                 & Apache 2.0 \\
        Gender\_guesser         & 0.4.0                 & GPLv3 \\
        Deep\_translator        & 1.11.4                & MIT \\
        Matplotlib              & 3.10.0                & PSF \\
        Seaborn                 & 0.13.2                & BSD \\
        Pandas                  & 2.2.2                 & GPLv3 \\
        Transformers            & 4.53.0                & Apache 2.0 \\
        Apify\_client           & 1.10.0                & Apache 2.0 \\
        Requests                & 2.32.4                & Apache 2.0 \\
        SSMS                    & 21.2.5                & Propietaria \\
        PowerBI                 & 2.140.1577.0          & Propietaria \\
        Draw.io                 & 27.0.6                & Apache 2.0 \\
        Docker                  & 27.4.0                & Apache 2.0 \\
        SonarQube               & 24.0.2                & GPLv3 \\
        Google Maps Scraper     & 1.0                   & MIT \\
        8icons                  & 2025                  & Icons8 UML \\
        httpx                   & 0.28.1                & BSD \\
        Dotenv                  & 1.1.1                 & BSD \\
		\bottomrule
	\end{tabularx}
	\caption{Tabla de herramientas y licencias}
\end{table}

\subsubsection{Propiedad intelectual}

El proyecto se puede dividir en varios productos finales con diferentes licencias:

\begin{table}[H]
	\centering
	\begin{tabularx}{\linewidth}{ p{0.7\columnwidth} p{0.3\columnwidth} }
		\toprule
		\textbf{Recurso}             &\textbf{Licencia}\\
        \toprule
		Dataset                      & CC BY-NC-SA \\
        Cuadro de mando              & CC BY-NC-SA \\
        Código fuente                & CC BY-NC-SA \\
        Documentación                & CC BY-NC-SA \\
		\bottomrule
	\end{tabularx}
	\caption{Tabla de productos y licencias}
\end{table}

Esto implica que el dataset, cuadro de mando y documentación pueden ser utilizados por cualquier persona siempre que se mantenga la misma licencia y se reconozca al autor original, pero no se permite su uso comercial.
El código fuente, sin embargo, se puede utilizar, modificar y distribuir libremente siempre que se mantenga la misma licencia y se reconozca al autor original.
\apendice{Especificación de Requisitos}

\section{Introducción}

Una muestra de cómo podría ser una tabla de casos de uso:

% Caso de Uso 1 -> Consultar Experimentos.
\begin{table}[p]
	\centering
	\begin{tabularx}{\linewidth}{ p{0.21\columnwidth} p{0.71\columnwidth} }
		\toprule
		\textbf{CU-1}    & \textbf{Ejemplo de caso de uso}\\
		\toprule
		\textbf{Versión}              & 1.0    \\
		\textbf{Autor}                & Alumno \\
		\textbf{Requisitos asociados} & RF-xx, RF-xx \\
		\textbf{Descripción}          & La descripción del CU \\
		\textbf{Precondición}         & Precondiciones (podría haber más de una) \\
		\textbf{Acciones}             &
		\begin{enumerate}
			\def\labelenumi{\arabic{enumi}.}
			\tightlist
			\item Pasos del CU
			\item Pasos del CU (añadir tantos como sean necesarios)
		\end{enumerate}\\
		\textbf{Postcondición}        & Postcondiciones (podría haber más de una) \\
		\textbf{Excepciones}          & Excepciones \\
		\textbf{Importancia}          & Alta o Media o Baja... \\
		\bottomrule
	\end{tabularx}
	\caption{CU-1 Nombre del caso de uso.}
\end{table}

\section{Objetivos generales}

\section{Catálogo de requisitos}

\section{Especificación de requisitos}



\apendice{Especificación de diseño}

\section{Introducción}

\section{Diseño de datos}

\section{Diseño arquitectónico}

\section{Diseño procedimental}




\apendice{Documentación técnica de programación}

\section{Introducción}

\section{Estructura de directorios}

\section{Manual del programador}

\section{Compilación, instalación y ejecución del proyecto}

\section{Pruebas del sistema}

\apendice{Documentación de usuario}

\section{Introducción}

En este apéndice se presenta la documentación de usuario necesaria para que el usuario final pueda utilizar el sistema desarrollado en este proyecto.
Este apéndice está enfocado al usuario final, que en principio no tiene conocimientos técnicos sobre el sistema y cuyo objetivo es usar el cuadro de mando final.

\section{Requisitos de usuarios}

El usuario final debe contar con varios requisitos para poder utilizar el cuadro de mando:
\begin{itemize}
    \item Tener acceso a un navegador web moderno con conexión a Internet.
    \item Tener acceso a una cuenta de Power BI y su aplicación descargada para poder acceder al cuadro de mando.
\end{itemize}

\section{Instalación}

La 'instalación' del cuadro de mando es un proceso muy simple. Simplemente hay que abrir el fichero \texttt{dashboardFinal.pbix} con la aplicación de Power BI Desktop. Una vez abierto, el usuario podrá ver el cuadro de mando y navegar por él.

\section{Manual del usuario}

El cuadro de mando está dividido en varias páginas. 
El usuario puede navegar de dos formas diferentes a través de dos componentes diferentes.

\subsection{Navegación por pestañas}

PowerBI cuenta con un menú inferior con las pestañas del cuadro de mando. Simplemente haciendo clic sobre ellas se puede navegar entre ellas.

\imagen{pestanas}{Pestañas de navegación}{1}

\subsection{Navegación por botones (sidebar)}

También se ha implementado de forma manual una barra lateral con iconos para navegar entre las difernetes pestañas. 
Para usarlos, el usuario debe pulsar el botón CONTROL a la vez que hace clic en el icono deseado.
Se muestra a continuación la barra lateral de forma apaisada. (En el cuadro de mando se encuentra a la izquierda en vertical).
\imagen{sidebar}{Barra lateral}{1}

\subsection{Filtros}

El cuadro de mando cuenta con dos filtros principales basados en los tipos de recursos y municipios. 
Estos filtros están disponibles en todas las pestañas.
\imagen{filtros}{Filtros de recursos y municipios}{0.5}

\subsection{Métricas}

Todas las gráficas del cuadro de mando son interactivas e intuitivas.
El usuario puede clicar en ellas para filtrar de la misma forma que con los filtros principales.
Las métricas principales del cuadro de mando son el TORI, el número de reseñas y la polaridad.
\imagen{metricas}{Métricas principales}{1}

\subsection{Borrado de filtros}

Se cuenta con un botón para borrar los filtros aplicados en el cuadro de mando.
Este botón se encuentra en la parte superior derecha y se acciona pulsando CONTROL a la vez que se hace clic sobre él.
\imagen{borrarFiltros}{Botón de borrado de filtros}{0.1}

\subsection{Ayuda}

Finalmente, cada pestaña cuenta con una sección de ayuda a la derecha que explica las funcionalidades desarrolladas en cada página.

\subsection{Páginas del cuadro de mando}
A continuación se muestran dos ejemplos de páginas del cuadro de mando en el que se ven los aspectos anteriores.
\imagen{recursos}{Página de recursos}{1}
\imagen{evolucionTemporal}{Página de evolución temporal}{1}
\apendice{Anexo de sostenibilización curricular}

\section{Introducción}

En este apéndice se explica la relación que existe entre el proyecto y los Objetivos de Desarrollo Sostenible (ODS) y las competencias de sostenibilidad adquiridas durante su desarrollo.
Para ello, se da una breve explicación de los ODS correspondientes y se relacionan con las tareas o funcionalidades realizadas en el proyecto.

Estos ODS fueron aprobados en 2015 por todos los miembros de la Organización de las Naciones Unidas (ONU) y tienen como objetivo 
acabar con la pobreza, mejorar la salud y la educación, reducir las desigualdades, garantizar la paz y prosperidad de la gente y del planeta e impulsar el crecimiento económico.\cite{ods}

\section{Objetivos de Desarrollo Sostenible}

El proyecto se centra en la sostenibilidad de los recursos turísticos de la provincia de Burgos, lo que implica una gestión responsable y sostenible de estos recursos.

\subsection{ODS 5: Igualdad de género}

Se proporcionan gráficas que permiten diferenciar las valoraciones y número de reseñas de los diferentes recursos dependiendo del género de los usuarios.
Esto permite identificar si hay alguna diferencia significativa en la valoración de los recursos turísticos entre hombres y mujeres, lo que puede ayudar a mejorar la oferta turística y adaptarla a las necesidades de ambos géneros.
También puede ayudar a identificar los recursos más valorados por un género u otro y que puede llevar a fomentar la visibilidad de los colectivos poco representados.

\subsection{ODS 9: Industria, innovación e infraestructura}

El desarrollo de este proyecto ha implicado el uso de tecnologías innovadoras como una red neuronal en relación al aprendizaje automático.
Además, la explotación de los datos obtenidos de forma pública permite mejorar la infraestructura turística de la provincia de Burgos.

\subsection{ODS 11: Ciudades y comunidades sostenibles}

Este proyecto contribuye a la sostenibilidad de las ciudades y comunidades al proporcionar información sobre los recursos de la provincia de Burgos.
A través del cuadro de mando, se puede identificar y analizar las zonas con un determinado tipo de recurso lo que permite detectar zonas mal valoradas o con pocos recursos que podrían ser mejoradas.

\section{Competencias de sostenibilidad adquiridas}

A través del desarrollo de este proyecto, he adquirido competencias de sostenibilidad definidas por la Conferencia de Rectores y Rectoras de las Universidades Españolas (CRUE)\cite{crue}.

\subsection{SOS1 - Competencia en la contextualización crítica del conocimiento estableciendo interrelaciones
con la problemática social, económica y ambiental, local y/o global}

Durante el desarrollo del proyecto se ha prestado atención al contexto social y territorial de la provincia de Burgos, 
entendiendo cómo la distribución de los puntos de interés (POIs) y la interacción de los ciudadanos con estos espacios a través de reseñas puede reflejar desigualdades económicas, concentración de servicios y posibles carencias estructurales. 

\subsection{SOS2 - Competencia en la utilización sostenible de recursos y en la prevención de impactos negativos sobre el medio natural y social}

El proyecto ha utilizado datos públicos de Google Maps, lo que implica un uso sostenible de los recursos digitales disponibles.
Además, en todo momento se ha tratado de optimizar y minimizar el impacto ambiental del proyecto, evitando el uso de recursos innecesarios y fomentando la reutilización de datos existentes.

\subsection{SOS4 - Competencia en la aplicación de principios éticos relacionados con los valores de la
sostenibilidad en los comportamientos personales y profesionales}

En relación a la ética, se ha procurado garantizar la transparencia en el uso de los datos y en la presentación de los resultados, evitando cualquier tipo de manipulación o tergiversación de la información.

\section{Conclusiones}

En conclusión, he aprendido a relacionar el proyecto con los ODS y a reflexionar sobre la sostenibilidad de los recursos turísticos de la provincia de Burgos.
He adquirido competencias de sostenibilidad que me permiten entender la importancia de la sostenibilidad en el desarrollo de proyectos y en la gestión de recursos turísticos.

\bibliographystyle{plain}
\bibliography{bibliografiaAnexos}

\end{document}